\subsection{Questão 21}

\subsubsection{Enunciado}

Em um circuito \(LC\) oscilante, \(L = 3{,}00\,\text{mH}\) e \(C = 2{,}70\,\mu\text{F}\). No instante \(t = 0\) a carga do capacitor é zero e a corrente é \(2{,}00\,\text{A}\).

\begin{enumerate}
    \item[(a)] Qual é a carga máxima do capacitor?
    
    \item[(b)] Em que instante de tempo \(t > 0\) a taxa com a qual a energia é armazenada no capacitor é máxima pela primeira vez?
    
    \item[(c)] Qual é o valor dessa taxa?
\end{enumerate}

\subsubsection{Resolução}


\vspace{1cm}
